% ============================================================
%  FICHIER ISOLÉ — Section 4.3.10 Captures d'écran
%  Smart Farm AI v3.0 — Mohamed Sayari
%
%  UTILISATION OVERLEAF / AUTRE ÉDITEUR CLOUD :
%    1) Téléversez ce fichier
%    2) Téléversez vos 15 captures dans un dossier "screenshots/"
%    3) Collez % ============================================================
%  FICHIER ISOLÉ — Section 4.3.10 Captures d'écran
%  Smart Farm AI v3.0 — Mohamed Sayari
%
%  UTILISATION OVERLEAF / AUTRE ÉDITEUR CLOUD :
%    1) Téléversez ce fichier
%    2) Téléversez vos 15 captures dans un dossier "screenshots/"
%    3) Collez % ============================================================
%  FICHIER ISOLÉ — Section 4.3.10 Captures d'écran
%  Smart Farm AI v3.0 — Mohamed Sayari
%
%  UTILISATION OVERLEAF / AUTRE ÉDITEUR CLOUD :
%    1) Téléversez ce fichier
%    2) Téléversez vos 15 captures dans un dossier "screenshots/"
%    3) Collez % ============================================================
%  FICHIER ISOLÉ — Section 4.3.10 Captures d'écran
%  Smart Farm AI v3.0 — Mohamed Sayari
%
%  UTILISATION OVERLEAF / AUTRE ÉDITEUR CLOUD :
%    1) Téléversez ce fichier
%    2) Téléversez vos 15 captures dans un dossier "screenshots/"
%    3) Collez \input{ch4_screenshots_only} dans votre document maître
%       (ou compilez de manière autonome avec le préambule ci-dessous)
%
%  COMPILATION AUTONOME (décommentez le bloc ci-dessous) :
% \documentclass[12pt,a4paper]{article}
% \usepackage[utf8]{inputenc}
% \usepackage[T1]{fontenc}
% \usepackage[french]{babel}
% \usepackage{graphicx}
% \graphicspath{{screenshots/}}
% \usepackage{caption,subcaption,float}
% \begin{document}
% ============================================================

% ── 4.3.10 Captures d'écran réelles de l'application ─────────────────────────
\subsection{Captures d'écran réelles de l'application}

Les captures d'écran suivantes illustrent l'interface réelle de Smart Farm
AI~v3.0 en cours d'exécution. Elles couvrent les neuf modules du tableau de
bord propriétaire (SPA React~18/Vite~5) et les six écrans principaux de
l'application mobile Worker PWA (Android, mode hors-ligne). Les images sont
organisées selon les deux profils d'accès de la plateforme.

% ─────────────────────────────────────────────────────────────────────────────
%  GROUPE 1 — Interface Propriétaire
% ─────────────────────────────────────────────────────────────────────────────
\subsubsection*{Interface Propriétaire}

% ── S01 : Landing Page ───────────────────────────────────────────────────────
\begin{figure}[H]
\centering
\includegraphics[width=0.88\linewidth]{screenshots/s01_landing.png}
\caption{Page d'accueil (Landing Page) --- point d'entrée unique avec accès
distinct Propriétaire (JWT HS256) et Ouvrier PWA (WhatsApp OTP).}
\label{fig:s01_landing}
\end{figure}

La page d'accueil constitue le point d'entrée unique de la plateforme.
Elle présente l'identité visuelle de Smart Farm AI~v3.0 et propose deux
boutons d'accès distincts~: l'interface propriétaire (tableau de bord complet
protégé par JWT HS256) et l'application mobile ouvrière (flux WhatsApp OTP
sans mot de passe). Cette page est servie statiquement par Nginx en production
et bénéficie d'un score Lighthouse PWA de~92/100, avec prise en charge du
protocole d'installation sur Android et iOS.

% ── S02 : Tableau de bord + détection incendie ───────────────────────────────
\begin{figure}[H]
\centering
\includegraphics[width=\linewidth]{screenshots/s02_dashboard.png}
\caption{Tableau de bord principal (Owner) --- KPIs temps réel, jauges IoT
ESP32, carte GIS et alerte de détection de fumée d'incendie par YOLOv11.}
\label{fig:s02_dashboard}
\end{figure}

Le tableau de bord agrège en une vue unique les indicateurs critiques de
l'exploitation~: nombre d'animaux actifs, alertes non lues, température
moyenne IoT et production journalière. Les jauges circulaires Recharts
affichent en temps réel l'humidité du sol (Nœud~A, DHT22) et le poids de
la ruche (Nœud~B, HX711). Un bandeau d'alerte rouge signale la détection de
fumée d'incendie par le modèle YOLOv11 déployé sur la caméra de surveillance,
avec localisation de la zone concernée sur la mini-carte GIS intégrée.
Les alertes sont diffusées simultanément par WebSocket et par notification
WhatsApp Business aux responsables de la ferme.

% ── S03 : Gestion du Bétail ──────────────────────────────────────────────────
\begin{figure}[H]
\centering
\includegraphics[width=\linewidth]{screenshots/s03_livestock.png}
\caption{Module Gestion du Bétail --- inventaire par espèce (6~catégories~:
bovins, ovins, caprins, volailles, lapins, abeilles) avec fiches détaillées
et score de santé individuel.}
\label{fig:s03_livestock}
\end{figure}

L'interface de gestion du bétail affiche six cartes d'espèces avec le nombre
de têtes actives, le score de santé moyen (calculé par moyenne pondérée des
dernières visites vétérinaires) et le statut d'alerte éventuel. Chaque carte
donne accès aux fiches individuelles, aux journaux de soins (vaccinations,
traitements médicamenteux), à l'historique de pesée et aux diagnostics IA
associés. Les données sont filtrables par ferme, par statut sanitaire et par
période temporelle. Le module supporte l'ajout d'animaux en lot via import CSV.

% ── S04 : Arbres & Plantations — YOLO Disease Detection ─────────────────────
\begin{figure}[H]
\centering
\includegraphics[width=\linewidth]{screenshots/s04_plants.png}
\caption{Module Arbres et Plantations --- détection de maladies foliaires
par YOLOv11 (mildiou, rouille, brûlure bactérienne) avec score de confiance
et recommandation Darija automatique.}
\label{fig:s04_plants}
\end{figure}

Le module végétal permet l'analyse phytosanitaire par upload d'image ou
capture directe depuis la caméra du dispositif. Le modèle YOLOv11 entraîné
sur un corpus de~3\,200 images annotées (mildiou, rouille, brûlure bactérienne,
déficit hydrique) affiche les boîtes englobantes colorées par classe et le
score de confiance associé. L'assistant Darija génère automatiquement une
recommandation de traitement en dialecte tunisien (Labess-7B via Ollama,
fallback Groq llama-3.3-70b), avec les sources RAG extraites de la base
ChromaDB (fiches AVFA, INRAT) affichées en regard du résultat.

% ── S05 : APICRAFT Apiculture — YOLO Bee Detection ───────────────────────────
\begin{figure}[H]
\centering
\includegraphics[width=\linewidth]{screenshots/s05_apiculture.png}
\caption{Module APICRAFT Apiculture --- détection YOLO des abeilles et
de la reine (35~ouvrières + 1~reine, confiance~95--97\%) sur image de cadre
de ruche en cours d'analyse.}
\label{fig:s05_apiculture}
\end{figure}

L'interface apicole APICRAFT intègre la détection visuelle des insectes par
le modèle yolo26n-obb entraîné sur~1\,847 images annotées (deux classes~:
\textit{bee} et \textit{queen}). Sur l'exemple présenté, le modèle identifie
35~ouvrières et~1~reine avec des scores de confiance de~95 à~97\%,
permettant une évaluation automatique de la densité de la colonie.
Le score COLOSS est mis à jour après chaque analyse, et une alerte est levée
si la reine n'est pas détectée dans les~72\,h suivant la dernière visite
documentée. L'historique des détections est persisté dans
\texttt{DiagnosticHistory} pour suivi longitudinal de la colonie.

% ── S06 : Entrepôt ───────────────────────────────────────────────────────────
\begin{figure}[H]
\centering
\includegraphics[width=\linewidth]{screenshots/s06_warehouse.png}
\caption{Module Entrepôt de la Ferme --- 10~catégories, 59~articles, gestion
des stocks avec alertes de seuil minimal configurable et assistant IA en arabe.}
\label{fig:s06_warehouse}
\end{figure}

Le gestionnaire d'entrepôt organise les approvisionnements en 10~catégories
(alimentation animale, médicaments vétérinaires, équipements apicoles,
intrants agricoles, outils, etc.) pour un total de~59~références actives.
Chaque article dispose d'un seuil d'alerte minimal configurable~: lorsque le
stock descend en dessous du seuil, une alerte apparaît dans le tableau de bord
et est diffusée via WebSocket en temps réel. Un assistant IA intégré accepte
les requêtes en arabe et en français pour rechercher des articles, vérifier
les niveaux de stock ou suggérer des commandes de réapprovisionnement basées
sur la consommation historique des 30~derniers jours.

% ── S07 : SovereignAssistant Darija ──────────────────────────────────────────
\begin{figure}[H]
\centering
\includegraphics[width=\linewidth]{screenshots/s07_assistant.png}
\caption{SovereignAssistant --- interface conversationnelle Darija avec
historique des sessions, sources RAG ChromaDB et métriques qualité
(BLEU-4\,=\,0.68, ROUGE-L\,=\,0.72) affichées en temps réel.}
\label{fig:s07_assistant}
\end{figure}

L'assistant souverain Darija offre une interface conversationnelle en dialecte
tunisien agricole~: le panneau gauche affiche l'historique des sessions
précédentes (persisté dans \texttt{DiagnosticHistory}), tandis que le panneau
droit présente la réponse générée par Labess-7B accompagnée des~3 sources
ChromaDB les plus pertinentes (top-3 HNSW, similarité cosinus). Les métriques
de qualité BLEU-4\,=\,0.68, ROUGE-L\,=\,0.72 et Cohen's~$\kappa$\,=\,0.76
sont affichées en bas de chaque réponse pour permettre à l'agriculteur
d'évaluer la fiabilité des conseils reçus. L'interface accepte également
des images (LLaVA-7B pour la description visuelle) et des documents PDF scannés
(OCR Tesseract \texttt{fra+ara+eng}).

% ── S08 : Rapports Stratégiques (Arabic RTL) ─────────────────────────────────
\begin{figure}[H]
\centering
\includegraphics[width=\linewidth]{screenshots/s08_reports.png}
\caption{Module Rapports Stratégiques --- moteur de génération IA en arabe (RTL)
avec exports PDF/Excel et historique des rapports classés par ferme et par type.}
\label{fig:s08_reports}
\end{figure}

Le module de rapports exploite le LLM pour générer automatiquement des analyses
stratégiques en arabe (mise en page RTL native via la propriété CSS
\texttt{direction: rtl} et \texttt{text-align: right}), en français ou en anglais.
Quatre types de rapports sont disponibles~: rapport mensuel zootechnique,
rapport sanitaire vétérinaire, analyse financière et bilan apicole saisonnier.
Chaque rapport embarque des graphiques Recharts, des tableaux de données et une
section de recommandations générée par Groq llama-3.3-70b. L'export PDF est
réalisé côté client via \texttt{html2canvas}~+~\texttt{jsPDF}, et l'export Excel
via \texttt{xlsx}. Les rapports sont sauvegardés en base et consultables depuis
l'historique filtrable par ferme et par période.

% ── S09 : Architecture / À Propos du Projet ──────────────────────────────────
\begin{figure}[H]
\centering
\includegraphics[width=\linewidth]{screenshots/s09_architecture.png}
\caption{Page Architecture (À Propos du Projet) --- diagramme interactif
Three.js des 9~modules, flux de données animés et technologies clés.}
\label{fig:s09_architecture}
\end{figure}

La page «~À Propos du Projet~» présente l'architecture technique complète
de Smart Farm AI~v3.0 sous forme d'un diagramme interactif Three.js~: les
neuf modules fonctionnels (Vision IA, Apiculture, Aviculture ERP, Entrepôt,
Cartographie GIS, Agent Darija, Télémétrie IoT, Rapports, PWA Worker) sont
représentés comme des nœuds orbitaux autour du cœur FastAPI, avec les flux de
données animés (WebSocket, REST, MQTT, RAG) rendus par un shader GLSL
personnalisé. Cette page sert à la fois de documentation visuelle pour les
développeurs et de vitrine technologique pour les parties prenantes du projet.

% ─────────────────────────────────────────────────────────────────────────────
%  GROUPE 2 — Application Mobile Worker PWA
% ─────────────────────────────────────────────────────────────────────────────
\subsubsection*{Application Mobile Worker PWA}

Les six écrans suivants illustrent l'application mobile ouvrière installée
sur Android en tant que PWA. L'interface adopte un thème sombre adapté aux
conditions lumineuses variables du terrain (champs, étables, ruchers), avec
des zones de frappe conformes aux recommandations WCAG~2.1
($\geq$~44\,px).

% ── S10 + S11 : Login + Home ─────────────────────────────────────────────────
\begin{figure}[H]
\centering
\begin{subfigure}[b]{0.48\textwidth}
  \centering
  \includegraphics[width=\linewidth]{screenshots/s10_worker_login.png}
  \caption{Connexion OTP WhatsApp --- numéro, envoi du code, validation}
  \label{fig:s10_worker_login}
\end{subfigure}
\hfill
\begin{subfigure}[b]{0.48\textwidth}
  \centering
  \includegraphics[width=\linewidth]{screenshots/s11_worker_home.png}
  \caption{Écran d'accueil ouvrier --- résumé tâches du jour, météo, sync}
  \label{fig:s11_worker_home}
\end{subfigure}
\caption{PWA Worker --- authentification WhatsApp OTP et écran d'accueil mobile.}
\label{fig:s10s11_worker}
\end{figure}

L'écran de connexion guide l'ouvrier en deux étapes~: saisie du numéro WhatsApp
enregistré par le propriétaire, puis validation du code OTP à~6 chiffres reçu
par message WhatsApp Business (TTL~5\,min). Aucun mot de passe n'est mémorisé
sur l'appareil~; le token JWT de~7~jours est stocké dans le \texttt{localStorage}
chiffré du Service Worker. L'écran d'accueil affiche un résumé des tâches
assignées pour la journée, la météo locale (OpenMeteo, sans clé API), le statut
de la connexion réseau (icône de synchronisation verte/orange) et la date de la
dernière synchronisation avec le backend FastAPI.

% ── S12 + S13 : Tâches + Signaler ────────────────────────────────────────────
\begin{figure}[H]
\centering
\begin{subfigure}[b]{0.48\textwidth}
  \centering
  \includegraphics[width=\linewidth]{screenshots/s12_worker_tasks.png}
  \caption{Liste des tâches prioritaires --- vaccination, alimentation, ruches}
  \label{fig:s12_worker_tasks}
\end{subfigure}
\hfill
\begin{subfigure}[b]{0.48\textwidth}
  \centering
  \includegraphics[width=\linewidth]{screenshots/s13_worker_report.png}
  \caption{Formulaire Signaler --- photo, géolocalisation GPS, notes textuelles}
  \label{fig:s13_worker_report}
\end{subfigure}
\caption{PWA Worker --- liste des tâches et formulaire de signalement terrain.}
\label{fig:s12s13_worker}
\end{figure}

La liste des tâches affiche les missions assignées par le propriétaire, triées
par priorité (URGENT en rouge, NORMAL en orange) et horodatées. Les tâches de
vaccination bovine et de collecte d'œufs apparaissent en priorité critique.
L'ouvrier peut marquer une tâche comme terminée d'un appui long ou via le bouton
de validation~; l'action est immédiatement enregistrée dans Dexie IndexedDB pour
synchronisation différée au rétablissement de la connexion. Le formulaire
\textit{Signaler} permet de joindre une photo (caméra ou galerie), d'associer
automatiquement la géolocalisation GPS et de saisir des notes textuelles libres,
entièrement fonctionnel hors-ligne grâce à la file d'attente IndexedDB.

% ── S14 + S15 : Paramètres + Scanner ─────────────────────────────────────────
\begin{figure}[H]
\centering
\begin{subfigure}[b]{0.48\textwidth}
  \centering
  \includegraphics[width=\linewidth]{screenshots/s14_worker_settings.png}
  \caption{Paramètres et synchronisation --- file Dexie, statut Service Worker}
  \label{fig:s14_worker_settings}
\end{subfigure}
\hfill
\begin{subfigure}[b]{0.48\textwidth}
  \centering
  \includegraphics[width=\linewidth]{screenshots/s15_worker_scanner.png}
  \caption{Scanner YOLO terrain --- sélection du modèle (bétail / ruche / plantes / incendie)}
  \label{fig:s15_worker_scanner}
\end{subfigure}
\caption{PWA Worker --- synchronisation hors-ligne et scanner YOLO de terrain.}
\label{fig:s14s15_worker}
\end{figure}

L'écran de paramètres affiche l'état de la file d'attente Dexie IndexedDB
(nombre d'opérations en attente, date de la dernière synchronisation réussie)
et permet de déclencher une synchronisation manuelle. Le statut du Service Worker
(actif/inactif) et la version du cache PWA installée sont visibles, avec un
bouton «~Mettre à jour~» qui recharge le Service Worker si une nouvelle version
est disponible via le backend. L'écran scanner donne accès aux quatre modèles
YOLO déployés sur l'appareil~: \textit{Bétail} (bovins, ovins, caprins),
\textit{Ruche} (yolo26n-obb, abeilles + reine), \textit{Plantes} (maladies
foliaires) et \textit{Incendie} (détection fumée/flammes). L'ouvrier sélectionne
le modèle adapté à sa situation, capture ou importe une image, et reçoit le
résultat d'inférence ($\leq$~180\,ms CPU) avec une recommandation Darija
générée localement par Labess-7B ou via Groq llama-3.3-70b en mode fallback.

% ============================================================
%  FIN DU FICHIER ISOLÉ
%  Pour la compilation autonome, décommentez : \end{document}
% \end{document}
% ============================================================
 dans votre document maître
%       (ou compilez de manière autonome avec le préambule ci-dessous)
%
%  COMPILATION AUTONOME (décommentez le bloc ci-dessous) :
% \documentclass[12pt,a4paper]{article}
% \usepackage[utf8]{inputenc}
% \usepackage[T1]{fontenc}
% \usepackage[french]{babel}
% \usepackage{graphicx}
% \graphicspath{{screenshots/}}
% \usepackage{caption,subcaption,float}
% \begin{document}
% ============================================================

% ── 4.3.10 Captures d'écran réelles de l'application ─────────────────────────
\subsection{Captures d'écran réelles de l'application}

Les captures d'écran suivantes illustrent l'interface réelle de Smart Farm
AI~v3.0 en cours d'exécution. Elles couvrent les neuf modules du tableau de
bord propriétaire (SPA React~18/Vite~5) et les six écrans principaux de
l'application mobile Worker PWA (Android, mode hors-ligne). Les images sont
organisées selon les deux profils d'accès de la plateforme.

% ─────────────────────────────────────────────────────────────────────────────
%  GROUPE 1 — Interface Propriétaire
% ─────────────────────────────────────────────────────────────────────────────
\subsubsection*{Interface Propriétaire}

% ── S01 : Landing Page ───────────────────────────────────────────────────────
\begin{figure}[H]
\centering
\includegraphics[width=0.88\linewidth]{screenshots/s01_landing.png}
\caption{Page d'accueil (Landing Page) --- point d'entrée unique avec accès
distinct Propriétaire (JWT HS256) et Ouvrier PWA (WhatsApp OTP).}
\label{fig:s01_landing}
\end{figure}

La page d'accueil constitue le point d'entrée unique de la plateforme.
Elle présente l'identité visuelle de Smart Farm AI~v3.0 et propose deux
boutons d'accès distincts~: l'interface propriétaire (tableau de bord complet
protégé par JWT HS256) et l'application mobile ouvrière (flux WhatsApp OTP
sans mot de passe). Cette page est servie statiquement par Nginx en production
et bénéficie d'un score Lighthouse PWA de~92/100, avec prise en charge du
protocole d'installation sur Android et iOS.

% ── S02 : Tableau de bord + détection incendie ───────────────────────────────
\begin{figure}[H]
\centering
\includegraphics[width=\linewidth]{screenshots/s02_dashboard.png}
\caption{Tableau de bord principal (Owner) --- KPIs temps réel, jauges IoT
ESP32, carte GIS et alerte de détection de fumée d'incendie par YOLOv11.}
\label{fig:s02_dashboard}
\end{figure}

Le tableau de bord agrège en une vue unique les indicateurs critiques de
l'exploitation~: nombre d'animaux actifs, alertes non lues, température
moyenne IoT et production journalière. Les jauges circulaires Recharts
affichent en temps réel l'humidité du sol (Nœud~A, DHT22) et le poids de
la ruche (Nœud~B, HX711). Un bandeau d'alerte rouge signale la détection de
fumée d'incendie par le modèle YOLOv11 déployé sur la caméra de surveillance,
avec localisation de la zone concernée sur la mini-carte GIS intégrée.
Les alertes sont diffusées simultanément par WebSocket et par notification
WhatsApp Business aux responsables de la ferme.

% ── S03 : Gestion du Bétail ──────────────────────────────────────────────────
\begin{figure}[H]
\centering
\includegraphics[width=\linewidth]{screenshots/s03_livestock.png}
\caption{Module Gestion du Bétail --- inventaire par espèce (6~catégories~:
bovins, ovins, caprins, volailles, lapins, abeilles) avec fiches détaillées
et score de santé individuel.}
\label{fig:s03_livestock}
\end{figure}

L'interface de gestion du bétail affiche six cartes d'espèces avec le nombre
de têtes actives, le score de santé moyen (calculé par moyenne pondérée des
dernières visites vétérinaires) et le statut d'alerte éventuel. Chaque carte
donne accès aux fiches individuelles, aux journaux de soins (vaccinations,
traitements médicamenteux), à l'historique de pesée et aux diagnostics IA
associés. Les données sont filtrables par ferme, par statut sanitaire et par
période temporelle. Le module supporte l'ajout d'animaux en lot via import CSV.

% ── S04 : Arbres & Plantations — YOLO Disease Detection ─────────────────────
\begin{figure}[H]
\centering
\includegraphics[width=\linewidth]{screenshots/s04_plants.png}
\caption{Module Arbres et Plantations --- détection de maladies foliaires
par YOLOv11 (mildiou, rouille, brûlure bactérienne) avec score de confiance
et recommandation Darija automatique.}
\label{fig:s04_plants}
\end{figure}

Le module végétal permet l'analyse phytosanitaire par upload d'image ou
capture directe depuis la caméra du dispositif. Le modèle YOLOv11 entraîné
sur un corpus de~3\,200 images annotées (mildiou, rouille, brûlure bactérienne,
déficit hydrique) affiche les boîtes englobantes colorées par classe et le
score de confiance associé. L'assistant Darija génère automatiquement une
recommandation de traitement en dialecte tunisien (Labess-7B via Ollama,
fallback Groq llama-3.3-70b), avec les sources RAG extraites de la base
ChromaDB (fiches AVFA, INRAT) affichées en regard du résultat.

% ── S05 : APICRAFT Apiculture — YOLO Bee Detection ───────────────────────────
\begin{figure}[H]
\centering
\includegraphics[width=\linewidth]{screenshots/s05_apiculture.png}
\caption{Module APICRAFT Apiculture --- détection YOLO des abeilles et
de la reine (35~ouvrières + 1~reine, confiance~95--97\%) sur image de cadre
de ruche en cours d'analyse.}
\label{fig:s05_apiculture}
\end{figure}

L'interface apicole APICRAFT intègre la détection visuelle des insectes par
le modèle yolo26n-obb entraîné sur~1\,847 images annotées (deux classes~:
\textit{bee} et \textit{queen}). Sur l'exemple présenté, le modèle identifie
35~ouvrières et~1~reine avec des scores de confiance de~95 à~97\%,
permettant une évaluation automatique de la densité de la colonie.
Le score COLOSS est mis à jour après chaque analyse, et une alerte est levée
si la reine n'est pas détectée dans les~72\,h suivant la dernière visite
documentée. L'historique des détections est persisté dans
\texttt{DiagnosticHistory} pour suivi longitudinal de la colonie.

% ── S06 : Entrepôt ───────────────────────────────────────────────────────────
\begin{figure}[H]
\centering
\includegraphics[width=\linewidth]{screenshots/s06_warehouse.png}
\caption{Module Entrepôt de la Ferme --- 10~catégories, 59~articles, gestion
des stocks avec alertes de seuil minimal configurable et assistant IA en arabe.}
\label{fig:s06_warehouse}
\end{figure}

Le gestionnaire d'entrepôt organise les approvisionnements en 10~catégories
(alimentation animale, médicaments vétérinaires, équipements apicoles,
intrants agricoles, outils, etc.) pour un total de~59~références actives.
Chaque article dispose d'un seuil d'alerte minimal configurable~: lorsque le
stock descend en dessous du seuil, une alerte apparaît dans le tableau de bord
et est diffusée via WebSocket en temps réel. Un assistant IA intégré accepte
les requêtes en arabe et en français pour rechercher des articles, vérifier
les niveaux de stock ou suggérer des commandes de réapprovisionnement basées
sur la consommation historique des 30~derniers jours.

% ── S07 : SovereignAssistant Darija ──────────────────────────────────────────
\begin{figure}[H]
\centering
\includegraphics[width=\linewidth]{screenshots/s07_assistant.png}
\caption{SovereignAssistant --- interface conversationnelle Darija avec
historique des sessions, sources RAG ChromaDB et métriques qualité
(BLEU-4\,=\,0.68, ROUGE-L\,=\,0.72) affichées en temps réel.}
\label{fig:s07_assistant}
\end{figure}

L'assistant souverain Darija offre une interface conversationnelle en dialecte
tunisien agricole~: le panneau gauche affiche l'historique des sessions
précédentes (persisté dans \texttt{DiagnosticHistory}), tandis que le panneau
droit présente la réponse générée par Labess-7B accompagnée des~3 sources
ChromaDB les plus pertinentes (top-3 HNSW, similarité cosinus). Les métriques
de qualité BLEU-4\,=\,0.68, ROUGE-L\,=\,0.72 et Cohen's~$\kappa$\,=\,0.76
sont affichées en bas de chaque réponse pour permettre à l'agriculteur
d'évaluer la fiabilité des conseils reçus. L'interface accepte également
des images (LLaVA-7B pour la description visuelle) et des documents PDF scannés
(OCR Tesseract \texttt{fra+ara+eng}).

% ── S08 : Rapports Stratégiques (Arabic RTL) ─────────────────────────────────
\begin{figure}[H]
\centering
\includegraphics[width=\linewidth]{screenshots/s08_reports.png}
\caption{Module Rapports Stratégiques --- moteur de génération IA en arabe (RTL)
avec exports PDF/Excel et historique des rapports classés par ferme et par type.}
\label{fig:s08_reports}
\end{figure}

Le module de rapports exploite le LLM pour générer automatiquement des analyses
stratégiques en arabe (mise en page RTL native via la propriété CSS
\texttt{direction: rtl} et \texttt{text-align: right}), en français ou en anglais.
Quatre types de rapports sont disponibles~: rapport mensuel zootechnique,
rapport sanitaire vétérinaire, analyse financière et bilan apicole saisonnier.
Chaque rapport embarque des graphiques Recharts, des tableaux de données et une
section de recommandations générée par Groq llama-3.3-70b. L'export PDF est
réalisé côté client via \texttt{html2canvas}~+~\texttt{jsPDF}, et l'export Excel
via \texttt{xlsx}. Les rapports sont sauvegardés en base et consultables depuis
l'historique filtrable par ferme et par période.

% ── S09 : Architecture / À Propos du Projet ──────────────────────────────────
\begin{figure}[H]
\centering
\includegraphics[width=\linewidth]{screenshots/s09_architecture.png}
\caption{Page Architecture (À Propos du Projet) --- diagramme interactif
Three.js des 9~modules, flux de données animés et technologies clés.}
\label{fig:s09_architecture}
\end{figure}

La page «~À Propos du Projet~» présente l'architecture technique complète
de Smart Farm AI~v3.0 sous forme d'un diagramme interactif Three.js~: les
neuf modules fonctionnels (Vision IA, Apiculture, Aviculture ERP, Entrepôt,
Cartographie GIS, Agent Darija, Télémétrie IoT, Rapports, PWA Worker) sont
représentés comme des nœuds orbitaux autour du cœur FastAPI, avec les flux de
données animés (WebSocket, REST, MQTT, RAG) rendus par un shader GLSL
personnalisé. Cette page sert à la fois de documentation visuelle pour les
développeurs et de vitrine technologique pour les parties prenantes du projet.

% ─────────────────────────────────────────────────────────────────────────────
%  GROUPE 2 — Application Mobile Worker PWA
% ─────────────────────────────────────────────────────────────────────────────
\subsubsection*{Application Mobile Worker PWA}

Les six écrans suivants illustrent l'application mobile ouvrière installée
sur Android en tant que PWA. L'interface adopte un thème sombre adapté aux
conditions lumineuses variables du terrain (champs, étables, ruchers), avec
des zones de frappe conformes aux recommandations WCAG~2.1
($\geq$~44\,px).

% ── S10 + S11 : Login + Home ─────────────────────────────────────────────────
\begin{figure}[H]
\centering
\begin{subfigure}[b]{0.48\textwidth}
  \centering
  \includegraphics[width=\linewidth]{screenshots/s10_worker_login.png}
  \caption{Connexion OTP WhatsApp --- numéro, envoi du code, validation}
  \label{fig:s10_worker_login}
\end{subfigure}
\hfill
\begin{subfigure}[b]{0.48\textwidth}
  \centering
  \includegraphics[width=\linewidth]{screenshots/s11_worker_home.png}
  \caption{Écran d'accueil ouvrier --- résumé tâches du jour, météo, sync}
  \label{fig:s11_worker_home}
\end{subfigure}
\caption{PWA Worker --- authentification WhatsApp OTP et écran d'accueil mobile.}
\label{fig:s10s11_worker}
\end{figure}

L'écran de connexion guide l'ouvrier en deux étapes~: saisie du numéro WhatsApp
enregistré par le propriétaire, puis validation du code OTP à~6 chiffres reçu
par message WhatsApp Business (TTL~5\,min). Aucun mot de passe n'est mémorisé
sur l'appareil~; le token JWT de~7~jours est stocké dans le \texttt{localStorage}
chiffré du Service Worker. L'écran d'accueil affiche un résumé des tâches
assignées pour la journée, la météo locale (OpenMeteo, sans clé API), le statut
de la connexion réseau (icône de synchronisation verte/orange) et la date de la
dernière synchronisation avec le backend FastAPI.

% ── S12 + S13 : Tâches + Signaler ────────────────────────────────────────────
\begin{figure}[H]
\centering
\begin{subfigure}[b]{0.48\textwidth}
  \centering
  \includegraphics[width=\linewidth]{screenshots/s12_worker_tasks.png}
  \caption{Liste des tâches prioritaires --- vaccination, alimentation, ruches}
  \label{fig:s12_worker_tasks}
\end{subfigure}
\hfill
\begin{subfigure}[b]{0.48\textwidth}
  \centering
  \includegraphics[width=\linewidth]{screenshots/s13_worker_report.png}
  \caption{Formulaire Signaler --- photo, géolocalisation GPS, notes textuelles}
  \label{fig:s13_worker_report}
\end{subfigure}
\caption{PWA Worker --- liste des tâches et formulaire de signalement terrain.}
\label{fig:s12s13_worker}
\end{figure}

La liste des tâches affiche les missions assignées par le propriétaire, triées
par priorité (URGENT en rouge, NORMAL en orange) et horodatées. Les tâches de
vaccination bovine et de collecte d'œufs apparaissent en priorité critique.
L'ouvrier peut marquer une tâche comme terminée d'un appui long ou via le bouton
de validation~; l'action est immédiatement enregistrée dans Dexie IndexedDB pour
synchronisation différée au rétablissement de la connexion. Le formulaire
\textit{Signaler} permet de joindre une photo (caméra ou galerie), d'associer
automatiquement la géolocalisation GPS et de saisir des notes textuelles libres,
entièrement fonctionnel hors-ligne grâce à la file d'attente IndexedDB.

% ── S14 + S15 : Paramètres + Scanner ─────────────────────────────────────────
\begin{figure}[H]
\centering
\begin{subfigure}[b]{0.48\textwidth}
  \centering
  \includegraphics[width=\linewidth]{screenshots/s14_worker_settings.png}
  \caption{Paramètres et synchronisation --- file Dexie, statut Service Worker}
  \label{fig:s14_worker_settings}
\end{subfigure}
\hfill
\begin{subfigure}[b]{0.48\textwidth}
  \centering
  \includegraphics[width=\linewidth]{screenshots/s15_worker_scanner.png}
  \caption{Scanner YOLO terrain --- sélection du modèle (bétail / ruche / plantes / incendie)}
  \label{fig:s15_worker_scanner}
\end{subfigure}
\caption{PWA Worker --- synchronisation hors-ligne et scanner YOLO de terrain.}
\label{fig:s14s15_worker}
\end{figure}

L'écran de paramètres affiche l'état de la file d'attente Dexie IndexedDB
(nombre d'opérations en attente, date de la dernière synchronisation réussie)
et permet de déclencher une synchronisation manuelle. Le statut du Service Worker
(actif/inactif) et la version du cache PWA installée sont visibles, avec un
bouton «~Mettre à jour~» qui recharge le Service Worker si une nouvelle version
est disponible via le backend. L'écran scanner donne accès aux quatre modèles
YOLO déployés sur l'appareil~: \textit{Bétail} (bovins, ovins, caprins),
\textit{Ruche} (yolo26n-obb, abeilles + reine), \textit{Plantes} (maladies
foliaires) et \textit{Incendie} (détection fumée/flammes). L'ouvrier sélectionne
le modèle adapté à sa situation, capture ou importe une image, et reçoit le
résultat d'inférence ($\leq$~180\,ms CPU) avec une recommandation Darija
générée localement par Labess-7B ou via Groq llama-3.3-70b en mode fallback.

% ============================================================
%  FIN DU FICHIER ISOLÉ
%  Pour la compilation autonome, décommentez : \end{document}
% \end{document}
% ============================================================
 dans votre document maître
%       (ou compilez de manière autonome avec le préambule ci-dessous)
%
%  COMPILATION AUTONOME (décommentez le bloc ci-dessous) :
% \documentclass[12pt,a4paper]{article}
% \usepackage[utf8]{inputenc}
% \usepackage[T1]{fontenc}
% \usepackage[french]{babel}
% \usepackage{graphicx}
% \graphicspath{{screenshots/}}
% \usepackage{caption,subcaption,float}
% \begin{document}
% ============================================================

% ── 4.3.10 Captures d'écran réelles de l'application ─────────────────────────
\subsection{Captures d'écran réelles de l'application}

Les captures d'écran suivantes illustrent l'interface réelle de Smart Farm
AI~v3.0 en cours d'exécution. Elles couvrent les neuf modules du tableau de
bord propriétaire (SPA React~18/Vite~5) et les six écrans principaux de
l'application mobile Worker PWA (Android, mode hors-ligne). Les images sont
organisées selon les deux profils d'accès de la plateforme.

% ─────────────────────────────────────────────────────────────────────────────
%  GROUPE 1 — Interface Propriétaire
% ─────────────────────────────────────────────────────────────────────────────
\subsubsection*{Interface Propriétaire}

% ── S01 : Landing Page ───────────────────────────────────────────────────────
\begin{figure}[H]
\centering
\includegraphics[width=0.88\linewidth]{screenshots/s01_landing.png}
\caption{Page d'accueil (Landing Page) --- point d'entrée unique avec accès
distinct Propriétaire (JWT HS256) et Ouvrier PWA (WhatsApp OTP).}
\label{fig:s01_landing}
\end{figure}

La page d'accueil constitue le point d'entrée unique de la plateforme.
Elle présente l'identité visuelle de Smart Farm AI~v3.0 et propose deux
boutons d'accès distincts~: l'interface propriétaire (tableau de bord complet
protégé par JWT HS256) et l'application mobile ouvrière (flux WhatsApp OTP
sans mot de passe). Cette page est servie statiquement par Nginx en production
et bénéficie d'un score Lighthouse PWA de~92/100, avec prise en charge du
protocole d'installation sur Android et iOS.

% ── S02 : Tableau de bord + détection incendie ───────────────────────────────
\begin{figure}[H]
\centering
\includegraphics[width=\linewidth]{screenshots/s02_dashboard.png}
\caption{Tableau de bord principal (Owner) --- KPIs temps réel, jauges IoT
ESP32, carte GIS et alerte de détection de fumée d'incendie par YOLOv11.}
\label{fig:s02_dashboard}
\end{figure}

Le tableau de bord agrège en une vue unique les indicateurs critiques de
l'exploitation~: nombre d'animaux actifs, alertes non lues, température
moyenne IoT et production journalière. Les jauges circulaires Recharts
affichent en temps réel l'humidité du sol (Nœud~A, DHT22) et le poids de
la ruche (Nœud~B, HX711). Un bandeau d'alerte rouge signale la détection de
fumée d'incendie par le modèle YOLOv11 déployé sur la caméra de surveillance,
avec localisation de la zone concernée sur la mini-carte GIS intégrée.
Les alertes sont diffusées simultanément par WebSocket et par notification
WhatsApp Business aux responsables de la ferme.

% ── S03 : Gestion du Bétail ──────────────────────────────────────────────────
\begin{figure}[H]
\centering
\includegraphics[width=\linewidth]{screenshots/s03_livestock.png}
\caption{Module Gestion du Bétail --- inventaire par espèce (6~catégories~:
bovins, ovins, caprins, volailles, lapins, abeilles) avec fiches détaillées
et score de santé individuel.}
\label{fig:s03_livestock}
\end{figure}

L'interface de gestion du bétail affiche six cartes d'espèces avec le nombre
de têtes actives, le score de santé moyen (calculé par moyenne pondérée des
dernières visites vétérinaires) et le statut d'alerte éventuel. Chaque carte
donne accès aux fiches individuelles, aux journaux de soins (vaccinations,
traitements médicamenteux), à l'historique de pesée et aux diagnostics IA
associés. Les données sont filtrables par ferme, par statut sanitaire et par
période temporelle. Le module supporte l'ajout d'animaux en lot via import CSV.

% ── S04 : Arbres & Plantations — YOLO Disease Detection ─────────────────────
\begin{figure}[H]
\centering
\includegraphics[width=\linewidth]{screenshots/s04_plants.png}
\caption{Module Arbres et Plantations --- détection de maladies foliaires
par YOLOv11 (mildiou, rouille, brûlure bactérienne) avec score de confiance
et recommandation Darija automatique.}
\label{fig:s04_plants}
\end{figure}

Le module végétal permet l'analyse phytosanitaire par upload d'image ou
capture directe depuis la caméra du dispositif. Le modèle YOLOv11 entraîné
sur un corpus de~3\,200 images annotées (mildiou, rouille, brûlure bactérienne,
déficit hydrique) affiche les boîtes englobantes colorées par classe et le
score de confiance associé. L'assistant Darija génère automatiquement une
recommandation de traitement en dialecte tunisien (Labess-7B via Ollama,
fallback Groq llama-3.3-70b), avec les sources RAG extraites de la base
ChromaDB (fiches AVFA, INRAT) affichées en regard du résultat.

% ── S05 : APICRAFT Apiculture — YOLO Bee Detection ───────────────────────────
\begin{figure}[H]
\centering
\includegraphics[width=\linewidth]{screenshots/s05_apiculture.png}
\caption{Module APICRAFT Apiculture --- détection YOLO des abeilles et
de la reine (35~ouvrières + 1~reine, confiance~95--97\%) sur image de cadre
de ruche en cours d'analyse.}
\label{fig:s05_apiculture}
\end{figure}

L'interface apicole APICRAFT intègre la détection visuelle des insectes par
le modèle yolo26n-obb entraîné sur~1\,847 images annotées (deux classes~:
\textit{bee} et \textit{queen}). Sur l'exemple présenté, le modèle identifie
35~ouvrières et~1~reine avec des scores de confiance de~95 à~97\%,
permettant une évaluation automatique de la densité de la colonie.
Le score COLOSS est mis à jour après chaque analyse, et une alerte est levée
si la reine n'est pas détectée dans les~72\,h suivant la dernière visite
documentée. L'historique des détections est persisté dans
\texttt{DiagnosticHistory} pour suivi longitudinal de la colonie.

% ── S06 : Entrepôt ───────────────────────────────────────────────────────────
\begin{figure}[H]
\centering
\includegraphics[width=\linewidth]{screenshots/s06_warehouse.png}
\caption{Module Entrepôt de la Ferme --- 10~catégories, 59~articles, gestion
des stocks avec alertes de seuil minimal configurable et assistant IA en arabe.}
\label{fig:s06_warehouse}
\end{figure}

Le gestionnaire d'entrepôt organise les approvisionnements en 10~catégories
(alimentation animale, médicaments vétérinaires, équipements apicoles,
intrants agricoles, outils, etc.) pour un total de~59~références actives.
Chaque article dispose d'un seuil d'alerte minimal configurable~: lorsque le
stock descend en dessous du seuil, une alerte apparaît dans le tableau de bord
et est diffusée via WebSocket en temps réel. Un assistant IA intégré accepte
les requêtes en arabe et en français pour rechercher des articles, vérifier
les niveaux de stock ou suggérer des commandes de réapprovisionnement basées
sur la consommation historique des 30~derniers jours.

% ── S07 : SovereignAssistant Darija ──────────────────────────────────────────
\begin{figure}[H]
\centering
\includegraphics[width=\linewidth]{screenshots/s07_assistant.png}
\caption{SovereignAssistant --- interface conversationnelle Darija avec
historique des sessions, sources RAG ChromaDB et métriques qualité
(BLEU-4\,=\,0.68, ROUGE-L\,=\,0.72) affichées en temps réel.}
\label{fig:s07_assistant}
\end{figure}

L'assistant souverain Darija offre une interface conversationnelle en dialecte
tunisien agricole~: le panneau gauche affiche l'historique des sessions
précédentes (persisté dans \texttt{DiagnosticHistory}), tandis que le panneau
droit présente la réponse générée par Labess-7B accompagnée des~3 sources
ChromaDB les plus pertinentes (top-3 HNSW, similarité cosinus). Les métriques
de qualité BLEU-4\,=\,0.68, ROUGE-L\,=\,0.72 et Cohen's~$\kappa$\,=\,0.76
sont affichées en bas de chaque réponse pour permettre à l'agriculteur
d'évaluer la fiabilité des conseils reçus. L'interface accepte également
des images (LLaVA-7B pour la description visuelle) et des documents PDF scannés
(OCR Tesseract \texttt{fra+ara+eng}).

% ── S08 : Rapports Stratégiques (Arabic RTL) ─────────────────────────────────
\begin{figure}[H]
\centering
\includegraphics[width=\linewidth]{screenshots/s08_reports.png}
\caption{Module Rapports Stratégiques --- moteur de génération IA en arabe (RTL)
avec exports PDF/Excel et historique des rapports classés par ferme et par type.}
\label{fig:s08_reports}
\end{figure}

Le module de rapports exploite le LLM pour générer automatiquement des analyses
stratégiques en arabe (mise en page RTL native via la propriété CSS
\texttt{direction: rtl} et \texttt{text-align: right}), en français ou en anglais.
Quatre types de rapports sont disponibles~: rapport mensuel zootechnique,
rapport sanitaire vétérinaire, analyse financière et bilan apicole saisonnier.
Chaque rapport embarque des graphiques Recharts, des tableaux de données et une
section de recommandations générée par Groq llama-3.3-70b. L'export PDF est
réalisé côté client via \texttt{html2canvas}~+~\texttt{jsPDF}, et l'export Excel
via \texttt{xlsx}. Les rapports sont sauvegardés en base et consultables depuis
l'historique filtrable par ferme et par période.

% ── S09 : Architecture / À Propos du Projet ──────────────────────────────────
\begin{figure}[H]
\centering
\includegraphics[width=\linewidth]{screenshots/s09_architecture.png}
\caption{Page Architecture (À Propos du Projet) --- diagramme interactif
Three.js des 9~modules, flux de données animés et technologies clés.}
\label{fig:s09_architecture}
\end{figure}

La page «~À Propos du Projet~» présente l'architecture technique complète
de Smart Farm AI~v3.0 sous forme d'un diagramme interactif Three.js~: les
neuf modules fonctionnels (Vision IA, Apiculture, Aviculture ERP, Entrepôt,
Cartographie GIS, Agent Darija, Télémétrie IoT, Rapports, PWA Worker) sont
représentés comme des nœuds orbitaux autour du cœur FastAPI, avec les flux de
données animés (WebSocket, REST, MQTT, RAG) rendus par un shader GLSL
personnalisé. Cette page sert à la fois de documentation visuelle pour les
développeurs et de vitrine technologique pour les parties prenantes du projet.

% ─────────────────────────────────────────────────────────────────────────────
%  GROUPE 2 — Application Mobile Worker PWA
% ─────────────────────────────────────────────────────────────────────────────
\subsubsection*{Application Mobile Worker PWA}

Les six écrans suivants illustrent l'application mobile ouvrière installée
sur Android en tant que PWA. L'interface adopte un thème sombre adapté aux
conditions lumineuses variables du terrain (champs, étables, ruchers), avec
des zones de frappe conformes aux recommandations WCAG~2.1
($\geq$~44\,px).

% ── S10 + S11 : Login + Home ─────────────────────────────────────────────────
\begin{figure}[H]
\centering
\begin{subfigure}[b]{0.48\textwidth}
  \centering
  \includegraphics[width=\linewidth]{screenshots/s10_worker_login.png}
  \caption{Connexion OTP WhatsApp --- numéro, envoi du code, validation}
  \label{fig:s10_worker_login}
\end{subfigure}
\hfill
\begin{subfigure}[b]{0.48\textwidth}
  \centering
  \includegraphics[width=\linewidth]{screenshots/s11_worker_home.png}
  \caption{Écran d'accueil ouvrier --- résumé tâches du jour, météo, sync}
  \label{fig:s11_worker_home}
\end{subfigure}
\caption{PWA Worker --- authentification WhatsApp OTP et écran d'accueil mobile.}
\label{fig:s10s11_worker}
\end{figure}

L'écran de connexion guide l'ouvrier en deux étapes~: saisie du numéro WhatsApp
enregistré par le propriétaire, puis validation du code OTP à~6 chiffres reçu
par message WhatsApp Business (TTL~5\,min). Aucun mot de passe n'est mémorisé
sur l'appareil~; le token JWT de~7~jours est stocké dans le \texttt{localStorage}
chiffré du Service Worker. L'écran d'accueil affiche un résumé des tâches
assignées pour la journée, la météo locale (OpenMeteo, sans clé API), le statut
de la connexion réseau (icône de synchronisation verte/orange) et la date de la
dernière synchronisation avec le backend FastAPI.

% ── S12 + S13 : Tâches + Signaler ────────────────────────────────────────────
\begin{figure}[H]
\centering
\begin{subfigure}[b]{0.48\textwidth}
  \centering
  \includegraphics[width=\linewidth]{screenshots/s12_worker_tasks.png}
  \caption{Liste des tâches prioritaires --- vaccination, alimentation, ruches}
  \label{fig:s12_worker_tasks}
\end{subfigure}
\hfill
\begin{subfigure}[b]{0.48\textwidth}
  \centering
  \includegraphics[width=\linewidth]{screenshots/s13_worker_report.png}
  \caption{Formulaire Signaler --- photo, géolocalisation GPS, notes textuelles}
  \label{fig:s13_worker_report}
\end{subfigure}
\caption{PWA Worker --- liste des tâches et formulaire de signalement terrain.}
\label{fig:s12s13_worker}
\end{figure}

La liste des tâches affiche les missions assignées par le propriétaire, triées
par priorité (URGENT en rouge, NORMAL en orange) et horodatées. Les tâches de
vaccination bovine et de collecte d'œufs apparaissent en priorité critique.
L'ouvrier peut marquer une tâche comme terminée d'un appui long ou via le bouton
de validation~; l'action est immédiatement enregistrée dans Dexie IndexedDB pour
synchronisation différée au rétablissement de la connexion. Le formulaire
\textit{Signaler} permet de joindre une photo (caméra ou galerie), d'associer
automatiquement la géolocalisation GPS et de saisir des notes textuelles libres,
entièrement fonctionnel hors-ligne grâce à la file d'attente IndexedDB.

% ── S14 + S15 : Paramètres + Scanner ─────────────────────────────────────────
\begin{figure}[H]
\centering
\begin{subfigure}[b]{0.48\textwidth}
  \centering
  \includegraphics[width=\linewidth]{screenshots/s14_worker_settings.png}
  \caption{Paramètres et synchronisation --- file Dexie, statut Service Worker}
  \label{fig:s14_worker_settings}
\end{subfigure}
\hfill
\begin{subfigure}[b]{0.48\textwidth}
  \centering
  \includegraphics[width=\linewidth]{screenshots/s15_worker_scanner.png}
  \caption{Scanner YOLO terrain --- sélection du modèle (bétail / ruche / plantes / incendie)}
  \label{fig:s15_worker_scanner}
\end{subfigure}
\caption{PWA Worker --- synchronisation hors-ligne et scanner YOLO de terrain.}
\label{fig:s14s15_worker}
\end{figure}

L'écran de paramètres affiche l'état de la file d'attente Dexie IndexedDB
(nombre d'opérations en attente, date de la dernière synchronisation réussie)
et permet de déclencher une synchronisation manuelle. Le statut du Service Worker
(actif/inactif) et la version du cache PWA installée sont visibles, avec un
bouton «~Mettre à jour~» qui recharge le Service Worker si une nouvelle version
est disponible via le backend. L'écran scanner donne accès aux quatre modèles
YOLO déployés sur l'appareil~: \textit{Bétail} (bovins, ovins, caprins),
\textit{Ruche} (yolo26n-obb, abeilles + reine), \textit{Plantes} (maladies
foliaires) et \textit{Incendie} (détection fumée/flammes). L'ouvrier sélectionne
le modèle adapté à sa situation, capture ou importe une image, et reçoit le
résultat d'inférence ($\leq$~180\,ms CPU) avec une recommandation Darija
générée localement par Labess-7B ou via Groq llama-3.3-70b en mode fallback.

% ============================================================
%  FIN DU FICHIER ISOLÉ
%  Pour la compilation autonome, décommentez : \end{document}
% \end{document}
% ============================================================
 dans votre document maître
%       (ou compilez de manière autonome avec le préambule ci-dessous)
%
%  COMPILATION AUTONOME (décommentez le bloc ci-dessous) :
% \documentclass[12pt,a4paper]{article}
% \usepackage[utf8]{inputenc}
% \usepackage[T1]{fontenc}
% \usepackage[french]{babel}
% \usepackage{graphicx}
% \graphicspath{{screenshots/}}
% \usepackage{caption,subcaption,float}
% \begin{document}
% ============================================================

% ── 4.3.10 Captures d'écran réelles de l'application ─────────────────────────
\subsection{Captures d'écran réelles de l'application}

Les captures d'écran suivantes illustrent l'interface réelle de Smart Farm
AI~v3.0 en cours d'exécution. Elles couvrent les neuf modules du tableau de
bord propriétaire (SPA React~18/Vite~5) et les six écrans principaux de
l'application mobile Worker PWA (Android, mode hors-ligne). Les images sont
organisées selon les deux profils d'accès de la plateforme.

% ─────────────────────────────────────────────────────────────────────────────
%  GROUPE 1 — Interface Propriétaire
% ─────────────────────────────────────────────────────────────────────────────
\subsubsection*{Interface Propriétaire}

% ── S01 : Landing Page ───────────────────────────────────────────────────────
\begin{figure}[H]
\centering
\includegraphics[width=0.88\linewidth]{screenshots/s01_landing.png}
\caption{Page d'accueil (Landing Page) --- point d'entrée unique avec accès
distinct Propriétaire (JWT HS256) et Ouvrier PWA (WhatsApp OTP).}
\label{fig:s01_landing}
\end{figure}

La page d'accueil constitue le point d'entrée unique de la plateforme.
Elle présente l'identité visuelle de Smart Farm AI~v3.0 et propose deux
boutons d'accès distincts~: l'interface propriétaire (tableau de bord complet
protégé par JWT HS256) et l'application mobile ouvrière (flux WhatsApp OTP
sans mot de passe). Cette page est servie statiquement par Nginx en production
et bénéficie d'un score Lighthouse PWA de~92/100, avec prise en charge du
protocole d'installation sur Android et iOS.

% ── S02 : Tableau de bord + détection incendie ───────────────────────────────
\begin{figure}[H]
\centering
\includegraphics[width=\linewidth]{screenshots/s02_dashboard.png}
\caption{Tableau de bord principal (Owner) --- KPIs temps réel, jauges IoT
ESP32, carte GIS et alerte de détection de fumée d'incendie par YOLOv11.}
\label{fig:s02_dashboard}
\end{figure}

Le tableau de bord agrège en une vue unique les indicateurs critiques de
l'exploitation~: nombre d'animaux actifs, alertes non lues, température
moyenne IoT et production journalière. Les jauges circulaires Recharts
affichent en temps réel l'humidité du sol (Nœud~A, DHT22) et le poids de
la ruche (Nœud~B, HX711). Un bandeau d'alerte rouge signale la détection de
fumée d'incendie par le modèle YOLOv11 déployé sur la caméra de surveillance,
avec localisation de la zone concernée sur la mini-carte GIS intégrée.
Les alertes sont diffusées simultanément par WebSocket et par notification
WhatsApp Business aux responsables de la ferme.

% ── S03 : Gestion du Bétail ──────────────────────────────────────────────────
\begin{figure}[H]
\centering
\includegraphics[width=\linewidth]{screenshots/s03_livestock.png}
\caption{Module Gestion du Bétail --- inventaire par espèce (6~catégories~:
bovins, ovins, caprins, volailles, lapins, abeilles) avec fiches détaillées
et score de santé individuel.}
\label{fig:s03_livestock}
\end{figure}

L'interface de gestion du bétail affiche six cartes d'espèces avec le nombre
de têtes actives, le score de santé moyen (calculé par moyenne pondérée des
dernières visites vétérinaires) et le statut d'alerte éventuel. Chaque carte
donne accès aux fiches individuelles, aux journaux de soins (vaccinations,
traitements médicamenteux), à l'historique de pesée et aux diagnostics IA
associés. Les données sont filtrables par ferme, par statut sanitaire et par
période temporelle. Le module supporte l'ajout d'animaux en lot via import CSV.

% ── S04 : Arbres & Plantations — YOLO Disease Detection ─────────────────────
\begin{figure}[H]
\centering
\includegraphics[width=\linewidth]{screenshots/s04_plants.png}
\caption{Module Arbres et Plantations --- détection de maladies foliaires
par YOLOv11 (mildiou, rouille, brûlure bactérienne) avec score de confiance
et recommandation Darija automatique.}
\label{fig:s04_plants}
\end{figure}

Le module végétal permet l'analyse phytosanitaire par upload d'image ou
capture directe depuis la caméra du dispositif. Le modèle YOLOv11 entraîné
sur un corpus de~3\,200 images annotées (mildiou, rouille, brûlure bactérienne,
déficit hydrique) affiche les boîtes englobantes colorées par classe et le
score de confiance associé. L'assistant Darija génère automatiquement une
recommandation de traitement en dialecte tunisien (Labess-7B via Ollama,
fallback Groq llama-3.3-70b), avec les sources RAG extraites de la base
ChromaDB (fiches AVFA, INRAT) affichées en regard du résultat.

% ── S05 : APICRAFT Apiculture — YOLO Bee Detection ───────────────────────────
\begin{figure}[H]
\centering
\includegraphics[width=\linewidth]{screenshots/s05_apiculture.png}
\caption{Module APICRAFT Apiculture --- détection YOLO des abeilles et
de la reine (35~ouvrières + 1~reine, confiance~95--97\%) sur image de cadre
de ruche en cours d'analyse.}
\label{fig:s05_apiculture}
\end{figure}

L'interface apicole APICRAFT intègre la détection visuelle des insectes par
le modèle yolo26n-obb entraîné sur~1\,847 images annotées (deux classes~:
\textit{bee} et \textit{queen}). Sur l'exemple présenté, le modèle identifie
35~ouvrières et~1~reine avec des scores de confiance de~95 à~97\%,
permettant une évaluation automatique de la densité de la colonie.
Le score COLOSS est mis à jour après chaque analyse, et une alerte est levée
si la reine n'est pas détectée dans les~72\,h suivant la dernière visite
documentée. L'historique des détections est persisté dans
\texttt{DiagnosticHistory} pour suivi longitudinal de la colonie.

% ── S06 : Entrepôt ───────────────────────────────────────────────────────────
\begin{figure}[H]
\centering
\includegraphics[width=\linewidth]{screenshots/s06_warehouse.png}
\caption{Module Entrepôt de la Ferme --- 10~catégories, 59~articles, gestion
des stocks avec alertes de seuil minimal configurable et assistant IA en arabe.}
\label{fig:s06_warehouse}
\end{figure}

Le gestionnaire d'entrepôt organise les approvisionnements en 10~catégories
(alimentation animale, médicaments vétérinaires, équipements apicoles,
intrants agricoles, outils, etc.) pour un total de~59~références actives.
Chaque article dispose d'un seuil d'alerte minimal configurable~: lorsque le
stock descend en dessous du seuil, une alerte apparaît dans le tableau de bord
et est diffusée via WebSocket en temps réel. Un assistant IA intégré accepte
les requêtes en arabe et en français pour rechercher des articles, vérifier
les niveaux de stock ou suggérer des commandes de réapprovisionnement basées
sur la consommation historique des 30~derniers jours.

% ── S07 : SovereignAssistant Darija ──────────────────────────────────────────
\begin{figure}[H]
\centering
\includegraphics[width=\linewidth]{screenshots/s07_assistant.png}
\caption{SovereignAssistant --- interface conversationnelle Darija avec
historique des sessions, sources RAG ChromaDB et métriques qualité
(BLEU-4\,=\,0.68, ROUGE-L\,=\,0.72) affichées en temps réel.}
\label{fig:s07_assistant}
\end{figure}

L'assistant souverain Darija offre une interface conversationnelle en dialecte
tunisien agricole~: le panneau gauche affiche l'historique des sessions
précédentes (persisté dans \texttt{DiagnosticHistory}), tandis que le panneau
droit présente la réponse générée par Labess-7B accompagnée des~3 sources
ChromaDB les plus pertinentes (top-3 HNSW, similarité cosinus). Les métriques
de qualité BLEU-4\,=\,0.68, ROUGE-L\,=\,0.72 et Cohen's~$\kappa$\,=\,0.76
sont affichées en bas de chaque réponse pour permettre à l'agriculteur
d'évaluer la fiabilité des conseils reçus. L'interface accepte également
des images (LLaVA-7B pour la description visuelle) et des documents PDF scannés
(OCR Tesseract \texttt{fra+ara+eng}).

% ── S08 : Rapports Stratégiques (Arabic RTL) ─────────────────────────────────
\begin{figure}[H]
\centering
\includegraphics[width=\linewidth]{screenshots/s08_reports.png}
\caption{Module Rapports Stratégiques --- moteur de génération IA en arabe (RTL)
avec exports PDF/Excel et historique des rapports classés par ferme et par type.}
\label{fig:s08_reports}
\end{figure}

Le module de rapports exploite le LLM pour générer automatiquement des analyses
stratégiques en arabe (mise en page RTL native via la propriété CSS
\texttt{direction: rtl} et \texttt{text-align: right}), en français ou en anglais.
Quatre types de rapports sont disponibles~: rapport mensuel zootechnique,
rapport sanitaire vétérinaire, analyse financière et bilan apicole saisonnier.
Chaque rapport embarque des graphiques Recharts, des tableaux de données et une
section de recommandations générée par Groq llama-3.3-70b. L'export PDF est
réalisé côté client via \texttt{html2canvas}~+~\texttt{jsPDF}, et l'export Excel
via \texttt{xlsx}. Les rapports sont sauvegardés en base et consultables depuis
l'historique filtrable par ferme et par période.

% ── S09 : Architecture / À Propos du Projet ──────────────────────────────────
\begin{figure}[H]
\centering
\includegraphics[width=\linewidth]{screenshots/s09_architecture.png}
\caption{Page Architecture (À Propos du Projet) --- diagramme interactif
Three.js des 9~modules, flux de données animés et technologies clés.}
\label{fig:s09_architecture}
\end{figure}

La page «~À Propos du Projet~» présente l'architecture technique complète
de Smart Farm AI~v3.0 sous forme d'un diagramme interactif Three.js~: les
neuf modules fonctionnels (Vision IA, Apiculture, Aviculture ERP, Entrepôt,
Cartographie GIS, Agent Darija, Télémétrie IoT, Rapports, PWA Worker) sont
représentés comme des nœuds orbitaux autour du cœur FastAPI, avec les flux de
données animés (WebSocket, REST, MQTT, RAG) rendus par un shader GLSL
personnalisé. Cette page sert à la fois de documentation visuelle pour les
développeurs et de vitrine technologique pour les parties prenantes du projet.

% ─────────────────────────────────────────────────────────────────────────────
%  GROUPE 2 — Application Mobile Worker PWA
% ─────────────────────────────────────────────────────────────────────────────
\subsubsection*{Application Mobile Worker PWA}

Les six écrans suivants illustrent l'application mobile ouvrière installée
sur Android en tant que PWA. L'interface adopte un thème sombre adapté aux
conditions lumineuses variables du terrain (champs, étables, ruchers), avec
des zones de frappe conformes aux recommandations WCAG~2.1
($\geq$~44\,px).

% ── S10 + S11 : Login + Home ─────────────────────────────────────────────────
\begin{figure}[H]
\centering
\begin{subfigure}[b]{0.48\textwidth}
  \centering
  \includegraphics[width=\linewidth]{screenshots/s10_worker_login.png}
  \caption{Connexion OTP WhatsApp --- numéro, envoi du code, validation}
  \label{fig:s10_worker_login}
\end{subfigure}
\hfill
\begin{subfigure}[b]{0.48\textwidth}
  \centering
  \includegraphics[width=\linewidth]{screenshots/s11_worker_home.png}
  \caption{Écran d'accueil ouvrier --- résumé tâches du jour, météo, sync}
  \label{fig:s11_worker_home}
\end{subfigure}
\caption{PWA Worker --- authentification WhatsApp OTP et écran d'accueil mobile.}
\label{fig:s10s11_worker}
\end{figure}

L'écran de connexion guide l'ouvrier en deux étapes~: saisie du numéro WhatsApp
enregistré par le propriétaire, puis validation du code OTP à~6 chiffres reçu
par message WhatsApp Business (TTL~5\,min). Aucun mot de passe n'est mémorisé
sur l'appareil~; le token JWT de~7~jours est stocké dans le \texttt{localStorage}
chiffré du Service Worker. L'écran d'accueil affiche un résumé des tâches
assignées pour la journée, la météo locale (OpenMeteo, sans clé API), le statut
de la connexion réseau (icône de synchronisation verte/orange) et la date de la
dernière synchronisation avec le backend FastAPI.

% ── S12 + S13 : Tâches + Signaler ────────────────────────────────────────────
\begin{figure}[H]
\centering
\begin{subfigure}[b]{0.48\textwidth}
  \centering
  \includegraphics[width=\linewidth]{screenshots/s12_worker_tasks.png}
  \caption{Liste des tâches prioritaires --- vaccination, alimentation, ruches}
  \label{fig:s12_worker_tasks}
\end{subfigure}
\hfill
\begin{subfigure}[b]{0.48\textwidth}
  \centering
  \includegraphics[width=\linewidth]{screenshots/s13_worker_report.png}
  \caption{Formulaire Signaler --- photo, géolocalisation GPS, notes textuelles}
  \label{fig:s13_worker_report}
\end{subfigure}
\caption{PWA Worker --- liste des tâches et formulaire de signalement terrain.}
\label{fig:s12s13_worker}
\end{figure}

La liste des tâches affiche les missions assignées par le propriétaire, triées
par priorité (URGENT en rouge, NORMAL en orange) et horodatées. Les tâches de
vaccination bovine et de collecte d'œufs apparaissent en priorité critique.
L'ouvrier peut marquer une tâche comme terminée d'un appui long ou via le bouton
de validation~; l'action est immédiatement enregistrée dans Dexie IndexedDB pour
synchronisation différée au rétablissement de la connexion. Le formulaire
\textit{Signaler} permet de joindre une photo (caméra ou galerie), d'associer
automatiquement la géolocalisation GPS et de saisir des notes textuelles libres,
entièrement fonctionnel hors-ligne grâce à la file d'attente IndexedDB.

% ── S14 + S15 : Paramètres + Scanner ─────────────────────────────────────────
\begin{figure}[H]
\centering
\begin{subfigure}[b]{0.48\textwidth}
  \centering
  \includegraphics[width=\linewidth]{screenshots/s14_worker_settings.png}
  \caption{Paramètres et synchronisation --- file Dexie, statut Service Worker}
  \label{fig:s14_worker_settings}
\end{subfigure}
\hfill
\begin{subfigure}[b]{0.48\textwidth}
  \centering
  \includegraphics[width=\linewidth]{screenshots/s15_worker_scanner.png}
  \caption{Scanner YOLO terrain --- sélection du modèle (bétail / ruche / plantes / incendie)}
  \label{fig:s15_worker_scanner}
\end{subfigure}
\caption{PWA Worker --- synchronisation hors-ligne et scanner YOLO de terrain.}
\label{fig:s14s15_worker}
\end{figure}

L'écran de paramètres affiche l'état de la file d'attente Dexie IndexedDB
(nombre d'opérations en attente, date de la dernière synchronisation réussie)
et permet de déclencher une synchronisation manuelle. Le statut du Service Worker
(actif/inactif) et la version du cache PWA installée sont visibles, avec un
bouton «~Mettre à jour~» qui recharge le Service Worker si une nouvelle version
est disponible via le backend. L'écran scanner donne accès aux quatre modèles
YOLO déployés sur l'appareil~: \textit{Bétail} (bovins, ovins, caprins),
\textit{Ruche} (yolo26n-obb, abeilles + reine), \textit{Plantes} (maladies
foliaires) et \textit{Incendie} (détection fumée/flammes). L'ouvrier sélectionne
le modèle adapté à sa situation, capture ou importe une image, et reçoit le
résultat d'inférence ($\leq$~180\,ms CPU) avec une recommandation Darija
générée localement par Labess-7B ou via Groq llama-3.3-70b en mode fallback.

% ============================================================
%  FIN DU FICHIER ISOLÉ
%  Pour la compilation autonome, décommentez : \end{document}
% \end{document}
% ============================================================
